\begin{question}
\noindent In a science lab, a student is modelling the path of a laser beam through a circular petri dish (representing a lens cross-section) for a Model UN science-diplomacy exhibit on optical technology.

\noindent A ray of light travels in a straight line along $TPS$, which is tangent to the circular dish at point $P$. Inside the dish, the light path is modelled by the chord $PQ$, and a second chord $QR$ is drawn to a point $R$ on the circle, so that triangle $PQR$ is formed inside the circle, with $R$ on the arc on the opposite side of $PQ$ from where the tangent touches.

\noindent The angle between the tangent $TPS$ and the chord $PQ$, angle $TPQ$, is $71^\circ$.

\begin{center}
\begin{tikzpicture}[scale=2.2]
    % Point of tangency
    \coordinate (P) at (-90:1);

    % circle
    \draw[name path=circ] (0,0) circle (1);

    % tangent line through P
    \draw (P) -- ($(P)!1.5!-90:(0,0)$) coordinate (S) node[below right]{$S$};
    \draw (P) -- ($(P)!1.5!90:(0,0)$) coordinate (T) node[above left]{$T$};

    % chord PQ - Q on circle, chosen so angle TPQ = 71 degrees roughly
    \path[name path=PQline] (P) -- ($(P)!2!-40:(0,0)$);
    \path[name intersections={of=PQline and circ, by={Qtemp,Q}}];
    \draw (P) -- (Q) node[above right]{$Q$};

    % point R on major arc, opposite side of PQ from T
    \coordinate (R) at (160:1);
    \draw (P) -- (R) node[left]{$R$};
    \draw (Q) -- (R);

    \node[below]{};
    \node[above left] at (P) {};
    \node at (P) [below] {$P$};

    \pic [draw, angle eccentricity=1.4, angle radius=0.6cm, "$71^\circ$"] {angle = T--P--Q};
    \pic [draw, angle eccentricity=1.5, angle radius=0.5cm, "$x$"] {angle = Q--R--P};

\end{tikzpicture}
\end{center}

\noindent Calculate the size of angle $PRQ$, marked $x$ on the diagram, which represents the deflection angle of the light path.

\vspace{5cm}

\begin{flushright}
\answerequnit{x}{$^\circ$}{4}
\end{flushright}
\end{question}